\documentclass[12pt, a4paper]{article}

% --- Pacchetti Scientifici e Formattazione ---
\usepackage{amsmath, amssymb, amsfonts, amsthm}
\usepackage{booktabs, graphicx}
\usepackage[document]{ragged2e}
\usepackage{xcolor}

\usepackage[utf8]{inputenc}
\usepackage[T1]{fontenc}
\usepackage[italian]{babel}
\usepackage{mathptmx}

\usepackage[includeheadfoot]{geometry}
\geometry{
    a4paper,
    top=3cm,
    bottom=3cm,
    left=3.5cm,
    right=3.5cm,
}

\usepackage{setspace}
\setstretch{1.3}
\usepackage[figuresright]{rotating}
\usepackage{float}

\title{Operazioni CRUD}
\author{Mattia Cavina}
\date{03 febbraio 2026}

\begin{document}
\begin{abstract}
    IL presente documento è volto a definire tutte le operazioni CRUD di base per la funzionalità del database lato software.
    Questa analisi è volta alla definizione dei metodi ed oggi da generare per l'iterazione del sofware con Data Base.
    Il documento verrà diviso in base ai ruoli e compiti tratte dalle interviste.
    L'ordine dei ruoli è in ordine descriscente, partendo dall'andimi fino al visualizzatore. Chi viene prima ha tutti le
    facolta di eseguire le operazioni di chi viene dopo, ma non viceversa.
\end{abstract}
\section{Operazioni Admin}
\subsection{Tabelle di lookup ``tipi''}
Le tabelle presenti nel database di lookup "tipi ( es:tipi\_prodotti, tipi\_proprietà ecc \dots) sono da considerarsi , a meno di acsi particolari tabella
di sola lettura, in quanto sono tabelle che contengono dati che non devono essere modificati o eliminati, ma solo visualizzati e utilizzati per le relazioni con altre tabelle.
La loro modifica non avrà una intrefaccia dedicata nel software, ma sarà gestita direttamente dal database tramite script SQL , garantendo così l'integrità dei dati e la coerenza del sistema.
Lato software verranno predisposte delle classi Enum per rappresentare i dati presenti in queste tabelle, facilitando così l'utilizzo ed evitando operazioni di join non necessarie.
Alle tabelle tipi si aggiungono anche altre tabelle che non verranno gestita da interfaccia come:
\begin{itemize}
    \item \textbf{Tabella lingue}
    \item \textbf{Tabella nazioni}
    \item \textbf{Tabella tabelle}
    \item \textbf{Tabella zone}
\end{itemize}
\subsection{Operazioni sugli utenti}
\begin{itemize}
    \item \textbf{Login:} Permette all'admin di accedere al sistema inserendo le proprie credenziali.
    \item \textbf{Creazione Utente:} Consente all'admin di creare nuovi account utente, specificando ruolo e permessi.
    \item \textbf{Modifica Utente:} Permette all'admin di aggiornare le informazioni degli utenti esistenti, come nome, email e ruolo.
    \item \textbf{Eliminazione Logica Utente:} Consente all'admin di disabilitare un account utente senza eliminarlo fisicamente dal database, mantenendo i dati per eventuali audit.
    \item \textbf{Visualizzazione Utenti:} Permette all'admin di visualizzare la lista degli utenti registrati, con possibilità di filtrare e ordinare.
\end{itemize}
\subsection{Operazioni sui prodotti}
\subsubsection{Tabella prodotti}
\begin{itemize}
    \item \textbf{Creazione Prodotto:} Consente all'admin di aggiungere nuovi prodotti al catalogo, specificando nome, descrizione, prezzo e disponibilità.
    \item \textbf{Modifica Prodotto:} Permette all'admin di aggiornare le informazioni dei prodotti esistenti, come prezzo, descrizione e disponibilità.
    \item \textbf{Eliminazione Logica Prodotto:} Consente all'admin di disabilitare un prodotto senza eliminarlo fisicamente dal database, mantenendo i dati per eventuali audit.
    \item \textbf{Visualizzazione Prodotti:} Permette all'admin di visualizzare la lista dei prodotti disponibili, con possibilità di filtrare e ordinare.
\end{itemize}
\subsubsection{Tabella proprietà prodotto}
\begin{itemize}
    \item \textbf{Aggiunta proprietà Prodotto:} Consente all'admin di aggiungere nuovi valori alle proprietà dei prodotti tra quelle disponibili.
    \item \textbf{Rimozione proprietà Prodotto:} Permette all'admin di rimuovere valori dalle proprietà dei prodotti tra quelle disponibili.
    \item \textbf{Modifica proprietà Prodotto:} Consente all'admin di modificare i valori delle proprietà dei prodotti tra quelle disponibili.
\end{itemize}
\subsubsection{Tabella proprietà}
\begin{itemize}
    \item \textbf{Creazione :} Consente all'admin di creare nuove proprietà per i prodotti, specificando nome e tipo di dato.
    \item \textbf{Modifica :} Permette all'admin di aggiornare le informazioni delle proprietà esistenti, come nome e tipo di dato.
    \item \textbf{Eliminazione Logica:} Consente all'admin di disabilitare una proprietà senza eliminarla fisicamente dal database, mantenendo i dati per eventuali audit.
    \item \textbf{Visualizzazione:} Permette all'admin di visualizzare la lista delle proprietà disponibili, con possibilità di filtrare e ordinare.
\end{itemize}
\subsubsection{Tabella proprieta valori}
\begin{itemize}
    \item \textbf{Creazione valori standard di una proprietà:} Consente all'admin di creare nuovi valori predefiniti per una proprietà, specificando il valore e la proprietà associata.
    \item \textbf{Modifica:} Permette all'admin di aggiornare le informazioni dei valori delle proprietà standard.
    \item \textbf{Eliminazione logica:} Consente all'admin di disabilitare un valore di una proprietà senza eliminarlo fisicamente dal database, mantenendo i dati per eventuali audit.
    \item \textbf{Visualizzazione valori predefinii di una proprietà:} Permette all'admin di visualizzare la lista dei valori predefiniti associati a una proprietà, con possibilità di filtrare e ordinare.
\end{itemize}
\subsubsection{Tabelle banche e banche\_sedi}
L'admin avrà facolta di gestire tutte le le tipologie di banche, acnhe quelle interne a differenza degli altri ruoli.
\begin{itemize}
    \item \textbf{Creazione banca:} Consente all'admin di aggiungere nuove banche al sistema, specificando nome, indirizzo e contatti.
    \item \textbf{Modifica banca:} Permette all'admin di aggiornare le informazioni delle banche esistenti, come nome, indirizzo e contatti.
    \item \textbf{Eliminazione logica banca:} Consente all'admin di disabilitare una banca senza eliminarla fisicamente dal database, mantenendo i dati per eventuali audit.
    \item \textbf{Visualizzazione banche:} Permette all'admin di visualizzare la lista delle banche disponibili, con possibilità di filtrare e ordinare.
\end{itemize}
\subsubsection{Tabelle condizioni e condizioni\_lingue}
\begin{itemize}
    \item \textbf{Creazione condizioni di vendita:} Consente all'admin di aggiungere nuove condizioni di vendita al sistema e relativa traduzione, mappata nella tabella condizioni\_lingue.
    \item \textbf{Modifica condizioni di vendita:} Permette all'admin di aggiornare le informazioni delle condizioni di vendita esistenti.
    \item \textbf{Eliminazione logica condizioni di vendita:} Consente all'admin di disabilitare una condizione di vendita senza eliminarla fisicamente dal database, mantenendo i dati per eventuali audit.
    \item \textbf{Visualizzazione condizioni di vendita:} Permette all'admin di visualizzare la lista delle condizioni di vendita disponibili, con possibilità di filtrare e ordinare.
\end{itemize}


\section{Operazioni Commerciali}
\subsection{Creazione e gestione delle offerte}
La creazione delle offerte è la parte fondamentale dell'intero programma.
Le operazioni che verranno elencate di seguito sono in ordine cronologico dalla generazione alla conferma d'ordine.
Si sottindende per semplicita che nelle operazio di seguito il commerciale
abbia la facoltà di creare, modificare e visualizzare.
\subsubsection{Creazione documento}
\begin{enumerate}
    \item \textbf{Creazione del primo documento:}
          \begin{itemize}
              \item Se il cliente o la destinazione non esiste \(\leftarrow\) \textbf{creazione cliente o destinazione} \(\leftarrow\) aggiunta al documento.
              \item Il commerciale può editare le informazioni editabili presenti nella tabella documento.
              \item Se il riferimento esterno non è presente tra quelli già presenti \(\leftarrow\) \textbf{creazione nuovo utente(senza accesso al software)} \(\leftarrow\) aggiunta al documento.
              \item Compillazione della sezione dedicata alla gestione delle lingue nei documenti \textbf{documenti\_lingue}, questo per avere la mappa delle traduzioni richieste dal cliente.
          \end{itemize}

    \item \textbf{Ricerca del prodotto:}
          L'inserimento di una riga, per guidare alla compilazione corretta, può avvenire in tre modi distinti:
          \begin{itemize}
              \item Inserimento manuale \(\leftarrow\) compilazione di tutti i campi necessari alla creazione di una riga documento.
              \item Inserimento tramite macroprodotto \(\leftarrow\) scelta del macroprodotto in base alla descrizione e proprietà \(\leftarrow\) inserimento automatico dei componenenti del macroprodotto \(\leftarrow\) eventuale modifica dei componenti.
              \item Creazione di un capitolo nella tabella \textbf{tipi\_capitoli} \(\leftarrow\) :
                    \begin{itemize}
                        \item Inserimento manuale \(\leftarrow\) compilazione di tutti i campi necessari alla creazione di una riga documento.
                        \item Inserimento tramite macroprodotto \(\leftarrow\) scelta del macroprodotto in base alla descrizione e proprietà \(\leftarrow\) inserimento automatico dei componenenti del macroprodotto \(\leftarrow\) eventuale modifica dei componenti.
                    \end{itemize}
          \end{itemize}
          Ogni prodotto o macroprodotto verrà poi cercato mediante filtri sulla descrizione e sulle proprietà, per facilitare la ricerca e l'inserimento.
    \item \textbf{Inserimento riga documento:}
          A sua volta l'inserimento di un componente comporta:
          \begin{itemize}
              \item Se il componente è standard ha già un prezzo associato, quindi l'inserimento è diretto.
                    In alternativa va inserito manualmente ed eventualmente calcolato da fogli esterni.
              \item Scelta delle scontische di riga da applicare al singolo componente.
              \item Modifica, se necessaria della descrizione.
              \item Inserimento dei dati tecnici richiesti dal cliente, se presenti, come portata.
                    segnalati come ricambi, quindi
              \item identificare se il componente è opzionale  o meno, in modo da poterlo calcolare diversamente nel totale dell'offerta.
          \end{itemize}
    \item \textbf{Gestione scontistiche:}
          Una volta completata l'offerta si possono editare le scontische e vedere il totale propostosi al cliente, in modo da poterlo modificare in base alle esigenze.
    \item \textbf{Calcolo e costificazione del montaggio:} Mediante la sezione dedicata al montaggio è possibile calcolare il costo del montaggio in base a dei parametri
          inseriti dal commerciale come il numero di tecnici necessari, il numero di giorni di montaggio e la difficoltà del montaggio, in modo da poterlo
          aggiungere al totale dell'offerta.
          Il tutto verrà gestito mediante le tabelle \textbf{montaggi} e\textbf{montaggi\_operatori}.
    \item{\textbf{Gestione di più versioni della stessa offerta:}}
          Per gestire la pluralità delle offerte è necessario che:
          \begin{itemize}
              \item Il commerciale possa revisionare una offerta senza che questa venga sovrascritta, in modo da mantenere la cronologia delle modifiche e delle versioni.
                    La revisione sarà gestita sulle righe del documento:
                    \begin{itemize}
                        \item Se la riga viene aggiornata essa viene duplicata ponendo la data di modifica vecchia nella nuova riga ed aggiornando
                              la data di modifica della riga vecchia, in modo da mantenere le relazioni.
                        \item Se la riga viene eliminata essa viene disabilitata, in modo da mantenere la cronologia delle modifiche e delle versioni.
                        \item Se il componente viene sostituito esso è come se il vecchio venisse eliminato e il nuovo inserito, si mantiene però il numero di posizione.
                    \end{itemize}
              \item Il commerciale deve avere la possbilità rispetto a vecchie versioni di riportarle valide in caso di accetazione di una vecchia versione.
          \end{itemize}
    \item \textbf{Conferma d'ordine:} Una volta che l'offerta è stata accettata dal cliente, il commerciale può confermare l'ordine, che comporta:
          \begin{itemize}
              \item Creazione del documento d'ordine come copia del documento d'offerta, con inserimento dell'id dell'offerta nel campo padre.
              \item Generazione del numero d'ordine, che sarà univoco e progressivo.
              \item Aggiornamento dello stato del documento d'offerta a ``accettata''.
              \item Duplicazione delle righe dell'offerta nell'ordine con eliminazione dei componenti opzionali non accettati dal cliente.
              \item Inserimento dei dati mancanti a livello di offerta.
                    Questi dati possono essere:
                    \begin{itemize}
                        \item Data di consegna prevista.
                        \item Modalità di pagamento.
                        \item Modalità di spedizione.
                        \item Eventuali note o istruzioni speciali per la produzione o la logistica.
                    \end{itemize}
          \end{itemize}

\end{enumerate}


\subsection{Tabella banche e banche\_sedi}
Il commerciale può gestire solo le banche esterne. ovvero quelle dei clienti.
\begin{itemize}
    \item \textbf{Creazione:} Consente al commerciale di aggiungere nuove banche esterne al sistema, specificando nome, indirizzo e contatti.
    \item \textbf{Modifica:} Permette al commerciale di aggiornare le informazioni delle banche esterne esistenti, come nome, indirizzo e contatti.
    \item \textbf{Eliminazione:} Consente al commerciale di disabilitare una banca esterna senza eliminarla fisicamente dal database, mantenendo i dati per eventuali audit.
    \item \textbf{Visualizzazione:} Permette al commerciale di visualizzare la lista delle banche esterne disponibili, con possibilità di filtrare e ordinare.
\end{itemize}

\section{Operazioni Project Manager}
Il project manager una volta che gli arriva da vagliare un ordine può:
\begin{itemize}
    \item Creazione del documento tecnico come copia del documento d'ordine, con inserimento dell'id dell'ordine nel campo padre.
    \item Duplicazione delle righe dell'ordine nel documento tecnico con eliminazione dei componenti opzionali non accettati dal cliente.
    \item Correzione dei componenti inseriti dal commerciale, in modo da adattarli alla realtà produttiva.
    \item Aggiunta o eliminazione di componenti, in modo da adattare l'offerta alla realtà produttiva.
    \item Correzione dei dati produttivi per la produzione, come tempi di consegna, modalità di spedizione e note per la produzione.
\end{itemize}

\section{Operazioni Logistica}
\subsection{Tabella corieri}
La logistica insieme all'admin ha facolta di gestire i corrieri, in quanto sono una risorsa fondamentale per la gestione delle spedizioni e dei resi.
Gli altri ruoli non avranno facoltà di getstire questa tabella.
\begin{itemize}
    \item \textbf{Creazione:} creazione di nuovi corieri.
    \item \textbf{Modifica:} modifica dei corieri.
    \item \textbf{Eliminazione:} eliminazione dei corieri
    \item \textbf{Visualizzazione:} visualizzazione dei corieri.
\end{itemize}

\section{Operazioni Visualizzatore}
Il visualizzatore ha facoltà di visualizzare tutti i documenti, ma non di modificarli, in quanto è un ruolo che serve solo per la consultazione dei dati.
Il visualizzatore deve poter vedere:
\begin{itemize}
    \item \textbf{Documenti:} visualizzazione di tutti i documenti, con possibilità di filtrare e ordinare.
    \item \textbf{Righe documento:} visualizzazione di tutte le righe dei documenti, con possibilità di filtrare e ordinare.
\end{itemize}
il visualizzatore può eseguire operazioni di inserimneto solo con le tabelle:
\begin{itemize}
    \item \textbf{Tabella feedback:} il visualizzatore può inserire segnalazioni o feeback sul programma.
    \item \textbf{Tabella notifiche:} il visualizzatore può inserire segnalazioni o notifiche per gli altri utenti.
\end{itemize}

\section{Mappatura oggetti nel software}
Definiti i ruoli e le operazioni è necessario mappare gli oggetti del database con le classi del software,
in modo da poter implementare le operazioni CRUD in modo efficiente e coerente con la struttura del database.
\subsection{Enum}
Le tabelle di lookup ``tipi'' verranno mappate con delle classi Enum, in modo da facilitare l'utilizzo e evitare operazioni di join non necessarie.
Le tabelle di lookup ``tipi'' sono:
\begin{itemize}
    \item \textbf{Tabella lingue}
    \item \textbf{Tabella nazioni}
    \item \textbf{Tabella tipi\_causale}
    \item \textbf{Tabella tipi\_cliente}
    \item \textbf{Tabella tipi\_contropartita}
    \item \textbf{Tabella tipi\_imballo}(da valutare)
    \item \textbf{Tabella tipi\_impianto}(da valutare)
    \item \textbf{Tabella tipi\_merceologia}(da valutare)
    \item \textbf{Tabella tipi\_operatore}
    \item \textbf{Tabella tipi\_resa}
    \item \textbf{Tabella tipi\_stato\_offerta}
    \item \textbf{Tabella tipi\_valuta} solo quelle attualmente utilizzate.
    \item \textbf{Elenco delle Tabelle} Elencare tutte le tabelle con relativo nome e descrizione,
          in modo da avere una panoramica completa delle tabelle presenti nel database.
    \item \
\end{itemize}
\subsection{Classi}
Le tabelle principali del database verranno mappate con delle classi, in modo da facilitare l'utilizzo e la gestione dei dati.
Le tabelle principali del database sono:
\begin{itemize}
    \item \textbf{Tabella banche}
    \item \textbf{Tabella banche\_sedi}
    \item \textbf{Tabella clienti}
    \item \textbf{Tabella clienti\_destinazioni}
    \item \textbf{Tabella condizioni}
    \item \textbf{Tabella condizioni\_lingue}
    \item \textbf{Tabella corrieri}
    \item \textbf{Tabella documenti}
    \item \textbf{Tabella documenti\_lingue}
    \item \textbf{Tabella documenti\_righe}
    \item \textbf{Tabella feedback}
    \item \textbf{Tabella montaggi}
    \item \textbf{Tabella montaggi\_operatori}
    \item \textbf{Tabella notifiche}
    \item \textbf{Tabella prodotti\_lingue}(Decidere se mappare le sue proprietà direttamente nella classe prodotto o meno)
    \item \textbf{Tabella prodotti\_proprieta}
    \item \textbf{Tabella proprieta}(Decidere se fare query solo per recuperare il nome della proprietà o mappare tutte le sue proprietà nella classe)
    \item \textbf{Tabella proprieta\_valori}
    \item \textbf{Tabella tipi\_pagamento}
    \item \textbf{Tabella tipi\_percentuale\_agente}
    \item \textbf{Tabella tipi\_proprieta}
    \item \textbf{Tabella tipi\_capitoli}(Decidere se fare query solo per recuperare il nome della proprietà o mappare tutte le sue proprietà nella classe)
    \item \textbf{Tabella utenti}
    \item \textbf{Tabella zone}
    \item \textbf{Tabella tabelle}
\end{itemize}

\subsection{Tabelle di relazione molti a molti}
In questa fase proviamo a non mappare le tabelle di relazione molti a molti,
in quanto sono tabelle che contengono solo le chiavi esterne delle tabelle principali, e non hanno una logica di business propria.
Questa tabelle sono:
\begin{itemize}
    \item \textbf{Tabella prodotti\_composizioni}
    \item \textbf{Tabella prodotti\_proprieta}
\end{itemize}
\end{document}